\begin{trapbox}
	\begin{description}[style=nextline, leftmargin=2.8em, labelsep=0.5em, itemsep=0.35em]
		\item[\textbf{\textcolor{CTrap}{易错点一（忽视前提条件）}}] 不等式中 $a$ 不一定为正数，若不注意 $a>0$ 直接套用等价形式，会出现方向错误。例如解 $|x|<-1$ 时写成 $-1<x<1$，实际解集为空。
		\item[\textbf{\textcolor{CTrap}{正确理解（确保前提）}}] 使用结论前必须检查 $a>0$；若 $a\le 0$，应结合绝对值非负性单独讨论。
		\item[\textbf{\textcolor{CTrap}{错因分析（忽略条件）}}] 对定理的前提条件记忆不牢，生搬硬套。
	\end{description}

	\medskip
	\begin{description}[style=nextline, leftmargin=2.8em, labelsep=0.5em, itemsep=0.35em]
		\item[\textbf{\textcolor{CTrap}{易错点二（混淆逻辑连接词）}}] 在 $|f(x)|>a$ 的等价转化中，误写为 $-a<f(x)<a$（即取“且”），导致解集错误。例如解 $|x|>2$ 错写成 $-2<x<2$。
		\item[\textbf{\textcolor{CTrap}{正确理解（分清两类）}}] $|f(x)|<a$ 对应“且”关系（交集），$|f(x)|>a$ 对应“或”关系（并集）。
		\item[\textbf{\textcolor{CTrap}{错因分析（记忆混淆）}}] 将两种情形的不等关系混淆，未区分中间的连接词。
	\end{description}

	\medskip
	\begin{description}[style=nextline, leftmargin=2.8em, labelsep=0.5em, itemsep=0.35em]
		\item[\textbf{\textcolor{CTrap}{易错点三（符号处理错误）}}] 解含绝对值不等式时，对 $f(x)$ 系数为负的情况处理不当。例如解 $|-2x+1|<3$，套用结论得 $-3<-2x+1<3$，但继续解时忘记除以负数改变不等号方向。
		\item[\textbf{\textcolor{CTrap}{正确理解（规范运算）}}] 套用结论后，正常解线性不等式，注意移项或除以负数时不等号反向。也可先处理系数为正的表达式（如 $|2x-1|<3$ 更方便）。
		\item[\textbf{\textcolor{CTrap}{错因分析（运算粗心）}}] 忽视常规不等式基本操作，只关注绝对值转化。
	\end{description}
\end{trapbox}
